\documentclass[10pt,letterpaper]{article}
\usepackage[margin=0.5in, top=0.4in, bottom=0.4in]{geometry}
\usepackage{enumitem}
\usepackage{titlesec}
\usepackage{hyperref}
\usepackage[T1]{fontenc}
\usepackage{xcolor}

% Remove page numbers
\pagestyle{empty}

% Section formatting
\titleformat{\section}{\normalsize\bfseries\uppercase}{}{0em}{}[\titlerule]
\titlespacing{\section}{0pt}{6pt}{4pt}

% Tighter lists
\setlist[itemize]{nosep, leftmargin=14pt, topsep=1pt, parsep=0pt, itemsep=1.5pt}

% Commands
\newcommand{\role}[4]{%
  \vspace{3pt}
  \noindent\textbf{#1} \hfill #2 \\
  \textit{#3} \hfill \textit{#4}
  \vspace{1pt}
}

\newcommand{\project}[3]{%
  \vspace{3pt}
  \noindent\textbf{#1} \hfill \textit{#2} \\
  {\small #3}
  \vspace{1pt}
}

\begin{document}

% Header
\begin{center}
  {\Large\bfseries VIBHOR KASHMIRA} \\[3pt]
  Detroit, MI \enspace$\cdot$\enspace (313) 746-0667 \enspace$\cdot$\enspace vibhorkashmira@gmail.com \enspace$\cdot$\enspace \href{https://bellerophon95.github.io/mypage}{bellerophon95.github.io/mypage}
\end{center}

\vspace{-4pt}
\noindent Software engineer with 7+ years building event-driven systems, data pipelines, and AI-powered products. Built high-throughput Kafka pipelines processing 10K+ msgs/sec at enterprise scale. Currently building backend systems for LLM agent workflows, context management, and evaluation tooling.

% Experience
\section{Experience}

\role{Software Engineer (AI)}{Aug 2025 -- Present}{Vizly Inc.}{Southfield, MI}
\begin{itemize}
  \item Architected distributed agent execution system with a full ReAct reasoning loop, human-in-the-loop approval gates, autonomous actions layer, and recursive fallback. Coordinated decoupled components via database-backed message passing, cron-driven scheduling, idempotent retries, and durable state persistence across human interrupts.
  \item Built MCP server (FastAPI) exposing tools for AI agents to load code samples, read/write to the database, and query metadata for autonomous task execution.
  \item Designed context management layer (Supabase) decoupling heavy datasets from LLM prompts, reducing token usage by 60\%.
  \item Architected RPyC bridge on AWS EKS for streaming S3 datasets to JupyterHub. Deployed with Docker/ECR, Prometheus, Grafana, and OpenTelemetry. Cut infrastructure costs by 10\% via FinOps.
\end{itemize}

\role{Software Engineer}{Jun 2017 -- Jun 2023}{Sprinklr Inc.}{Gurugram, India}
\begin{itemize}
  \item Led migration of legacy Engagement feed APIs to GraphQL, reducing response time by 15\% and network payload by 25\%. Leveraged gRPC and reactive processing (Spring WebFlux) for key performance paths.
  \item Engineered event-driven architecture using Kafka, processing 10,000+ messages/sec with end-to-end observability and event tracing.
  \item Optimized MongoDB via progressive sharding driven by event tracing data, reducing query latency by 15\% for the Engagement feed and 14 Governance asset classes.
  \item Reduced infrastructure costs by 10\% via AWS Lambda/EC2 for skewed workloads. Achieved 99.9\% availability across Kubernetes environments with near-zero-downtime deployments and automated testing (Gatling, REST Assured) via Jenkins CI/CD.
  \item Mentored junior developers through code reviews and pair programming, reducing regression issues by 20\%.
\end{itemize}

% Projects
\section{Projects}

\project{Rentaroost}{Jul -- Aug 2024}{\href{https://bellerophon95.github.io/mypage/rentaroost/}{bellerophon95.github.io/mypage/rentaroost/}}
\begin{itemize}
  \item Built real-time dynamic pricing pipeline using Apache Kafka and Apache Flink with watermark handling for out-of-order events. Prices adjust based on view and booking event streams with configurable TTLs via Redis.
  \item Implemented distributed transaction orchestration using the Kafka Saga pattern across payment (Stripe) and booking services with Spring WebFlux for non-blocking write flows.
  \item Designed CQRS-style architecture with GraphQL (Netflix DGS) for user-facing queries, gRPC for synchronous internal reads, and Kafka for async mutations. Deployed on Kubernetes.
\end{itemize}

\project{Nexus -- Multi-Agent RAG System}{Mar -- Apr 2026}{\href{https://project-nexus.duckdns.org}{project-nexus.duckdns.org} $\cdot$ \href{https://github.com/bellerophon95/nexus}{github.com/bellerophon95/nexus}}
\begin{itemize}
  \item Built and deployed a multi-agent RAG platform from zero to production on AWS in 8 days (143 commits, 3 hosting pivots). Retrieval pipeline combines Qdrant dense + Supabase BM25 sparse search with Cohere Rerank, feeding a LangGraph agent loop with Self-RAG hallucination gating. Integrated Langfuse tracing; real-time SSE streaming of agent state to Next.js 15 frontend. Runs on a \$15/month t3.small via Terraform IaC.
\end{itemize}

% Skills
\section{Technical Skills}

\noindent\textbf{Backend:} Python (FastAPI), Java (Spring Boot), gRPC, GraphQL (Netflix DGS), Spring WebFlux, REST \\
\textbf{Data \& Streaming:} Apache Kafka, Apache Flink, MongoDB, Redis, Supabase (PostgreSQL), Qdrant \\
\textbf{Infrastructure:} AWS (EKS, S3, Lambda, EC2), Kubernetes, Docker, Terraform, CI/CD (Jenkins, GitHub Actions) \\
\textbf{Observability:} Prometheus, Grafana, OpenTelemetry, ELK Stack, Langfuse \\
\textbf{AI Tooling:} LangGraph, RAG (hybrid retrieval, cross-encoder reranking), RAGAS, MCP, Vercel AI SDK

% Education
\section{Education}

\noindent\textbf{Wayne State University} -- M.S. in Computer Science (Dean's List) \hfill Aug 2023 -- Dec 2024 \\
\textbf{The LNMIIT} -- B.Tech in Computer Science \hfill Aug 2013 -- May 2017

\end{document}
